\documentclass[10pt,a4paper]{article}
\usepackage[T1]{fontenc}
\usepackage[margin=2cm]{geometry}
\usepackage{isabelle,isabellesym}
\usepackage{amsmath,amssymb}
\usepackage{pdfsetup}

\urlstyle{rm}
\isabellestyle{it}

\begin{document}

\title{Young's Inequality for Cosine Sums}
\author{Carles Mar\'in}
\maketitle

\begin{abstract}
W.\,H.~Young proved in 1913 that the partial sums of the cosine series $\sum_{k \geq 1} \cos(kx)/k$ are bounded
below by $-1$~\cite{Young1913}. We formalize this inequality: for every $n$ and every real $x$,
$\sum_{k=1}^{n} \cos(kx)/k \geq -1$. The proof takes a global minimum of the sum on $[0, 2\pi]$ and, at an
interior critical point, uses either the induction hypothesis or a double Abel summation with convex
coefficients.
\end{abstract}

\paragraph{Acknowledgement.} The formalization and the text of this document were prepared with the
assistance of generative AI systems, and checked by the author.

\tableofcontents

\section*{Introduction}

Write $c_n(x) = \sum_{k=1}^{n} \cos(kx)/k$. The bound is attained: $c_1(\pi) = -1$. Young's inequality is
the starting point of a series of sharper results on positive trigonometric sums; Brown and Koumandos sharpened
and unified several of them~\cite{BrownKoumandos1997}. The proof formalized here is by induction on $n$ and goes
in three steps.
\begin{enumerate}
\item By periodicity it is enough to bound $c_n$ on $[0, 2\pi]$, where it attains a minimum at some $s$. At
  $s = 0$ and $s = 2\pi$ the sum is nonnegative.
\item At an interior minimum the derivative $-\sum_{k=1}^{n} \sin(ks)$ vanishes. Multiplying by $2\sin(s/2)$
  and telescoping gives $\sin(ns/2)\,\sin((n+1)s/2) = 0$.
\item If $\sin(ns/2) = 0$, then $\cos(ns) = 1$ and $c_n(s) = c_{n-1}(s) + 1/n \geq -1$ by induction. If
  $\sin((n+1)s/2) = 0$, write $1 + c_n(s) = \sum_{k=0}^{n} a_k u_k$ with $a_0 = 2$, $a_k = 1/k$, $u_0 = 1/2$ and
  $u_k = \cos(ks)$, and sum by parts twice. The second differences of $a$ are nonnegative, the second partial
  sums of $u$ are $(1 - \cos((k+1)s)) / (4 \sin^2(s/2)) \geq 0$, and at the critical point the boundary terms
  are explicit, which gives $1 + c_n(s) \geq 0$.
\end{enumerate}

\parindent 0pt\parskip 0.5ex

\input{session}

\bibliographystyle{abbrv}
\bibliography{root}

\end{document}
